%% 2-7 何回もかける：累乗と指数（第11回で書き足し。ページ数は自由・偶数）
%% この章で定義した部品（頭文字 bkiih）：
%%   \bkiihlayers{(左下)}{枚数}{はば}{色}            重なった紙を横から見た図（1枚 0.9 mm）
%%   \bkiihcube{(左下)}{n}{1目もり mm}{色}            n×n×n の小さな立方体でできた立方体（見取図）
%%   \bkiihearth{(中心)}{半径}  \bkiihmoon{(中心)}{半径}   地球と月
%%   \bkiihswitch{(左下)}{0/1}                        スイッチ（1つ）
%%   \bkiihcell{(中心)}{半径}                          細菌（だ円）
\makeatletter
\newcommand\bkiihlayers[4]{\begin{scope}[shift={#1}]
  \foreach \k in {1,...,#2} {\fill[#4!45] (0,{(\k-1)*0.9}) rectangle (#3,{(\k-1)*0.9+0.6});}
  \end{scope}}
\newcommand\bkiihcube[4]{\begin{scope}[shift={#1}]
  \pgfmathsetmacro\bkiih@s{#2*#3}\pgfmathsetmacro\bkiih@dx{0.5*#2*#3}\pgfmathsetmacro\bkiih@dy{0.32*#2*#3}%
  \pgfmathtruncatemacro\bkiih@m{#2-1}%
  \filldraw[fill=#4!30, draw=#4!70!black, line width=0.7pt] (0,0) rectangle (\bkiih@s,\bkiih@s);
  \filldraw[fill=#4!12, draw=#4!70!black, line width=0.7pt] (0,\bkiih@s) -- (\bkiih@s,\bkiih@s) -- ++(\bkiih@dx,\bkiih@dy) -- (\bkiih@dx,\bkiih@s+\bkiih@dy) -- cycle;
  \filldraw[fill=#4!55, draw=#4!70!black, line width=0.7pt] (\bkiih@s,0) -- ++(\bkiih@dx,\bkiih@dy) -- ++(0,\bkiih@s) -- (\bkiih@s,\bkiih@s) -- cycle;
  \ifnum\bkiih@m>0
  \foreach \k in {1,...,\bkiih@m} {
    \draw[#4!70!black, line width=0.35pt] (\k*#3,0) -- (\k*#3,\bkiih@s) -- ++(\bkiih@dx,\bkiih@dy);
    \draw[#4!70!black, line width=0.35pt] (0,\k*#3) -- (\bkiih@s,\k*#3) -- ++(\bkiih@dx,\bkiih@dy);
    \draw[#4!70!black, line width=0.35pt] ({\k*\bkiih@dx/#2},{\bkiih@s+\k*\bkiih@dy/#2}) -- ++(\bkiih@s,0) -- ++(0,-\bkiih@s);}
  \fi
  \end{scope}}
\newcommand\bkiihearth[2]{\begin{scope}[shift={#1}]
  \fill[bkBlue!70] (0,0) circle (#2);
  \fill[bkGreen!70] (-0.45*#2,0.2*#2) .. controls (-0.1*#2,0.75*#2) and (0.3*#2,0.5*#2) .. (0.2*#2,0.05*#2)
     .. controls (0.1*#2,-0.4*#2) and (-0.5*#2,-0.3*#2) .. cycle;
  \fill[bkGreen!70] (0.45*#2,-0.35*#2) circle (0.18*#2);
  \draw[bkBlue!60!black, line width=0.6pt] (0,0) circle (#2);\end{scope}}
\newcommand\bkiihmoon[2]{\begin{scope}[shift={#1}]
  \fill[bkGray!35] (0,0) circle (#2);
  \fill[bkGray!55] (-0.3*#2,0.25*#2) circle (0.22*#2) (0.35*#2,-0.2*#2) circle (0.15*#2);
  \draw[bkGray!80!black, line width=0.6pt] (0,0) circle (#2);\end{scope}}
\newcommand\bkiihswitch[2]{\begin{scope}[shift={#1}]
  \filldraw[fill=white, draw=bkInk!70, line width=0.6pt, rounded corners=1mm] (0,0) rectangle (7,11);
  \ifnum#2=1 \fill[bkOrange, rounded corners=0.8mm] (1.2,5.5) rectangle (5.8,9.8);
  \else \fill[bkGray!45, rounded corners=0.8mm] (1.2,1.2) rectangle (5.8,5.5);\fi
  \node[font=\sffamily\bfseries\footnotesize, text=bkInk] at (3.5,-2.6) {#2};\end{scope}}
% 小さな見出し  \bkiihhead{見出し}
\newcommand\bkiihhead[1]{\par\addvspace{5pt}\noindent{\sffamily\bfseries\color{bkTheme!85!black}▶\ #1}\par\nopagebreak\vspace{1pt}}
\newcommand\bkiihcell[2]{\filldraw[fill=bkGreen!35, draw=bkGreen!70!black, line width=0.6pt] #1 ellipse ({#2} and {0.55*#2});
  \fill[bkGreen!70!black] #1 circle (0.18*#2);}
\makeatother

\bkchapter{2}{7}{\R{何回}{なん|かい}もかける：\\\R{累乗}{るい|じょう}と\R{指数}{し|すう}}{\bklevelrange{1}{2}}{\R{何回}{なん|かい}もかける：\R{累乗}{るい|じょう}と\R{指数}{し|すう}}
\bkmichishirube{1-2,1-21,2-1,2-2}{2-8,2-10,2-12,2-17,5-14}

\begin{goal}
\begin{itemize}[label={\color{bkTheme}●}, leftmargin=1.6\zw, itemsep=0pt, topsep=1pt]
\item \bklevel{1}\enspace \R{同}{おな}じ\R{数}{かず}を\R{何回}{なん|かい}もかけた\R{数}{かず}を、$2^3$ のような\R{累乗}{るい|じょう}で\R{書}{か}いて\R{計算}{けい|さん}できる。
\item \bklevel{1}\enspace \R{体積}{たい|せき}の\R{単位}{たん|い}（\unit{cm^3}・\unit{m^3}・L）や\R{大}{おお}きな\R{数}{かず}を、10の\R{累乗}{るい|じょう}で\R{表}{あらわ}せる。
\item \bklevel{2}\enspace \R{指数}{し|すう}のきまりで\R{計算}{けい|さん}でき、0\R{乗}{じょう}やマイナスの\R{指数}{し|すう}の\R{意味}{い|み}がわかる。
\end{itemize}
\end{goal}

\begin{hatena}
\noindent\begin{minipage}[c]{0.52\linewidth}
コピー\R{用紙}{よう|し}1\R{枚}{まい}の\R{厚}{あつ}さは、\R{約}{やく}\qty{0.1}{mm}です。\R{半分}{はん|ぶん}に\R{折}{お}ると、\R{厚}{あつ}さは2\R{倍}{ばい}。もう1\R{回}{かい}\R{折}{お}ると、そのまた2\R{倍}{ばい}になります。

もし\R{何回}{なん|かい}でも\R{折}{お}れるとしたら、42\R{回}{かい}\R{折}{お}った\R{紙}{かみ}の\R{厚}{あつ}さは、\R{地球}{ち|きゅう}から\R{月}{つき}（\R{約}{やく}38\R{万}{まん}\,km）までとどくでしょうか？
\end{minipage}\hfill
\begin{minipage}[c]{0.45\linewidth}\centering
\begin{tikzpicture}[x=1mm, y=1mm]
  \bkiihearth{(0,0)}{7}
  \bkiihmoon{(44,0)}{3.4}
  \fill[bkOrange!80] (7.6,-0.8) rectangle (33,0.8);
  \node[font=\sffamily\bfseries\Large, text=bkOrange!80!black] at (37.2,0) {？};
  \draw[bkGray, line width=0.6pt, <->] (0,-10) -- (44,-10);
  \node[anchor=north, font=\sffamily\footnotesize, text=bkInk] at (22,-11.5) {\R{約}{やく}38\R{万}{まん}\,km};
  \node[anchor=south, font=\sffamily\footnotesize, text=bkInk] at (0,8.5) {\R{地球}{ち|きゅう}};
  \node[anchor=south, font=\sffamily\footnotesize, text=bkInk] at (44,5) {\R{月}{つき}};
  \node[anchor=south, font=\sffamily\footnotesize, text=bkOrange!80!black] at (20.5,2.2) {42\R{回}{かい}\R{折}{お}った\R{紙}{かみ}};
\end{tikzpicture}
\end{minipage}
\end{hatena}

\R{折}{お}るたびに2\R{倍}{ばい}、2\R{倍}{ばい}、……。\R{同}{おな}じ\R{数}{かず}を\R{何回}{なん|かい}もかけると、\R{数}{かず}はおどろくほど\R{速}{はや}く\R{大}{おお}きくなります。この\R{章}{しょう}では、そんな\R{計算}{けい|さん}を\R{短}{みじか}く\R{書}{か}く\R{方法}{ほう|ほう}と、その\R{計算}{けい|さん}のきまりを\R{見}{み}つけます。

\clearpage
%% ---------------- p.89 やってみよう・思い出そう ----------------
\begin{yattemiyou}[\R{紙}{かみ}を\R{折}{お}って、\R{重}{かさ}なりを\R{数}{かぞ}えよう]
\noindent\begin{minipage}[c]{0.55\linewidth}
\begin{enumerate}[label={\bkcircled[bkGreen]{\arabic*}}, leftmargin=2.4\zw, itemsep=1pt, topsep=0pt]
\item コピー\R{用紙}{よう|し}を1\R{枚}{まい}\R{用意}{よう|い}して、\R{半分}{はん|ぶん}に\R{折}{お}る。
\item \R{折}{お}るたびに、\R{紙}{かみ}が\R{何枚}{なん|まい}\R{重}{かさ}なったかを\R{数}{かぞ}えて、ノートの\R{表}{ひょう}に\R{書}{か}こう（\R{折}{お}り\R{目}{め}の\R{反対}{はん|たい}がわから\R{数}{かぞ}えるとよい）。
\item 5\R{回}{かい}\R{折}{お}ったら、ものさしで\R{厚}{あつ}さをはかろう。
\item \R{何回}{なん|かい}まで\R{折}{お}れたかな？
\end{enumerate}
\end{minipage}\hfill
\begin{minipage}[c]{0.42\linewidth}\centering
\begin{tikzpicture}[x=1mm, y=1mm]
  \foreach \n/\h/\x in {0/1/0,1/2/12,2/4/24,3/8/36} {
    \bkiihlayers{(\x,0)}{\h}{9}{bkGreen}
    \node[anchor=north, font=\sffamily\footnotesize, text=bkInk] at (\x+4.5,-1.5) {\n\R{回}{かい}};
    \node[anchor=south, font=\sffamily\footnotesize, text=bkGreen!60!black] at (\x+4.5,8) {\h\R{枚}{まい}};}
\end{tikzpicture}\par\vspace{1pt}
{\sffamily\footnotesize \R{横}{よこ}から\R{見}{み}た\R{紙}{かみ}の\R{重}{かさ}なり}
\end{minipage}
\par\vspace{6pt}
\begin{bkcenter}
\begin{bktable}{colspec={Q[l,m,wd=31mm] *{7}{Q[c,m,wd=10.5mm]}}, rowsep=2pt, colsep=3pt, row{1}={bg=bkGreen!14}, column{1}={font=\sffamily\small}}
\R{折}{お}った\R{回数}{かい|すう} & 0 & 1 & 2 & 3 & 4 & 5 & 6 \\
\R{重}{かさ}なった\R{枚数}{まい|すう} & 1 & 2 & 4 & ？ & ？ & ？ & ？ \\
\R{厚}{あつ}さ（mm） & 0.1 & 0.2 & 0.4 & ？ & ？ & ？ & ？ \\
\end{bktable}
\end{bkcenter}
1\R{回}{かい}\R{折}{お}るごとに、\R{枚数}{まい|すう}は\R{何倍}{なん|ばい}になるかな？　5\R{回}{かい}\R{折}{お}った\R{厚}{あつ}さは、\R{計算}{けい|さん}すると\R{何}{なん}mmかな？
\end{yattemiyou}

\begin{omoidasou}
\bkrecall{1-2}{10が10こで100、100が10こで1000。10こ\R{集}{あつ}まると、1つ\R{左}{ひだり}の\R{位}{くらい}に\R{進}{すす}む。}
\bkrecall{1-21}{1\R{辺}{ぺん}が\qty{1}{cm}の\R{正方形}{せい|ほう|けい}の\R{広}{ひろ}さが\qty{1}{cm^2}。\R{長方形}{ちょう|ほう|けい}の\R{面積}{めん|せき}は、たて×よこ。}
\bkrecall{1-9}{$\qty{1}{L}=\qty{1000}{mL}$}
\bkrecall{2-1}{$\dfrac{1}{10}$にすると、\R{小数点}{しょう|すう|てん}が\R{左}{ひだり}へ1けた\R{動}{うご}く（1→0.1→0.01）。}
\end{omoidasou}

\clearpage
%% ---------------- 数学の道具 ----------------
\bksection{doguu}{\R{数学}{すう|がく}の\R{道具}{どう|ぐ}}

\bksubsection{1}{\R{同}{おな}じ\R{数}{かず}を\R{何回}{なん|かい}もかける：\R{累乗}{るい|じょう}}

\leavevmode\label{bk:2-7:ruijo}\bkyattemiyouref で、\R{紙}{かみ}を\R{折}{お}るたびに、\R{重}{かさ}なった\R{枚数}{まい|すう}は2\R{倍}{ばい}になりました。\R{枚数}{まい|すう}を、かけ\R{算}{ざん}の\R{式}{しき}で\R{書}{か}いてみましょう。

\begin{bkfigure}
\begin{tikzpicture}[x=1mm, y=1mm]
  \foreach \n/\h/\x/\e in {0/1/0/1, 1/2/27/2, 2/4/54/2\times2, 3/8/81/2\times2\times2, 4/16/108/2\times2\times2\times2} {
    \bkiihlayers{(\x,0)}{\h}{14}{bkTheme}
    \node[anchor=north, font=\sffamily\small, text=bkInk] at (\x+7,-1.6) {\n\R{回}{かい}\R{折}{お}る};
    \node[anchor=north, font=\sffamily\small, text=bkTheme!80!black] at (\x+7,-8.6) {\h\R{枚}{まい}};
    \node[anchor=north, font=\small, text=bkInk] at (\x+7,-15) {$\e$};}
  \foreach \x in {0,27,54,81} \draw[bk link] (\x+15.5,3) -- node[above, bk link label, pos=0.45] {×2} (\x+25.5,3);
\end{tikzpicture}
\end{bkfigure}

\noindent\begin{minipage}[c]{0.6\linewidth}
3\R{回}{かい}\R{折}{お}ったときの $2\times2\times2$ は、2を3\R{回}{かい}かけた\R{数}{かず}です。これを\R{短}{みじか}く $2^3$ と\R{書}{か}き、「2の3\R{乗}{じょう}」と\R{読}{よ}みます。
\R{右}{みぎ}かたの\R{小}{ちい}さな\R{数}{かず}を\textbf{\R{指数}{し|すう}}といいます。\bkidx{しすう}{\R{指数}{し|すう}}
\end{minipage}\hfill
\begin{minipage}[c]{0.37\linewidth}\centering
\begin{tikzpicture}[x=1mm, y=1mm]
  \node[font=\fontsize{30}{32}\selectfont, text=bkInk, inner sep=0pt] (b) at (0,0) {$2$};
  \node[font=\fontsize{17}{18}\selectfont, text=bkTheme, inner sep=0pt, anchor=south west] (e) at ([xshift=0.6mm, yshift=-2mm]b.north east) {$3$};
  \node[anchor=west, font=\sffamily\footnotesize, text=bkInk, align=left, inner sep=0pt] (t) at ([xshift=9mm]e.east) {\R{指数}{し|すう}：2を\\\R{何回}{なん|かい}かけるか};
  \draw[bkTheme, line width=0.7pt, -{Stealth[length=4pt]}] ([xshift=-1.5mm]t.west) -- ([xshift=1.5mm]e.east);
  \node[anchor=north, font=\small, text=bkInk] at (12,-6) {$2^3=2\times2\times2=8$};
\end{tikzpicture}
\end{minipage}

$2^3$ のように、\R{同}{おな}じ\R{数}{かず}をいくつかかけたものを\textbf{\R{累乗}{るい|じょう}}といいます。\bkidx{るいじょう}{\R{累乗}{るい|じょう}}\R{折}{お}った\R{回数}{かい|すう}が、そのまま\R{指数}{し|すう}になっていますね。4\R{回}{かい}\R{折}{お}ると $2^4=16$\R{枚}{まい}です。1\R{回}{かい}だけのときは、2を1つだけ\R{書}{か}くので $2^1=2$ と\R{書}{か}きます。

\bkiihhead{\R{表}{ひょう}にすると、1つ\R{右}{みぎ}へ\R{行}{い}くごとに×2}
$2^4$ は、$2^3=8$ にもう1つ2をかけたものです。だから、\R{前}{まえ}の\R{答}{こた}えに2をかければ、\R{次}{つぎ}の\R{答}{こた}えが\R{出}{で}ます。2を\R{何回}{なん|かい}も\R{書}{か}かなくてすみます。

\begin{bkcenter}
\begin{bktable}{colspec={Q[l,m,wd=17mm] *{10}{Q[c,m,wd=9.6mm]}}, rowsep=1.5pt, colsep=1.5pt, cells={font=\small}, column{1}={font=\sffamily\small}, row{1}={bg=bkTheme!10}}
\R{累乗}{るい|じょう} & $2^1$ & $2^2$ & $2^3$ & $2^4$ & $2^5$ & $2^6$ & $2^7$ & $2^8$ & $2^9$ & $2^{10}$ \\
\R{答}{こた}え & 2 & 4 & 8 & 16 & 32 & 64 & 128 & 256 & 512 & 1024 \\
\end{bktable}
\bktabcaption{2の\R{累乗}{るい|じょう}。\R{右}{みぎ}へ1つ\R{進}{すす}むと、2を1\R{回}{かい}\R{多}{おお}くかけるので2\R{倍}{ばい}になる。}
\end{bkcenter}

\R{指数}{し|すう}がたった10なのに、\R{答}{こた}えは1000をこえました。$2^{10}=1024$ は、この\R{章}{しょう}で\R{何度}{なん|ど}も\R{使}{つか}います。

\bkiihhead{2\R{以外}{い|がい}の\R{数}{かず}でも\R{同}{おな}じ}
\R{累乗}{るい|じょう}は、どんな\R{数}{かず}でも\R{作}{つく}れます。3を2\R{回}{かい}かけた $3\times3=9$ は $3^2$（3の2\R{乗}{じょう}）、10を3\R{回}{かい}かけた $10\times10\times10=1000$ は $10^3$（10の3\R{乗}{じょう}）です。1は\R{何回}{なん|かい}かけても1なので、$1^5=1\times1\times1\times1\times1=1$ です。

\bksubsection{2}{\R{面積}{めん|せき}は2\R{乗}{じょう}、\R{体積}{たい|せき}は3\R{乗}{じょう}}

\noindent\begin{minipage}[c]{0.63\linewidth}
\hspace*{1\zw}1\R{辺}{ぺん}\qty{3}{cm}の\R{正方形}{せい|ほう|けい}の\R{面積}{めん|せき}は、$3\times3=3^2=9$ で\qty{9}{cm^2}です（1-21）。\unit{cm^2}の\R{小}{ちい}さな2も、「cmを2\R{回}{かい}かけた」というしるしです。

\hspace*{1\zw}では、1\R{辺}{ぺん}\qty{3}{cm}の\R{立方体}{りっ|ぽう|たい}（さいころの\R{形}{かたち}）の\R{大}{おお}きさは？　1\R{辺}{ぺん}\qty{1}{cm}の\R{立方体}{りっ|ぽう|たい}が\R{何}{なん}こ\R{入}{はい}るか\R{数}{かぞ}えると、1だんに $3\times3=9$こ、それが3だんで $3\times3\times3=3^3=27$こです。
\end{minipage}\hfill
\begin{minipage}[c]{0.34\linewidth}\centering
\begin{tikzpicture}[x=1mm, y=1mm]
  \foreach \i in {0,1,2} \foreach \j in {0,1,2} \filldraw[fill=bkBlue!25, draw=bkBlue!70!black, line width=0.5pt] (\i*5,\j*5) rectangle ++(5,5);
  \node[anchor=north, font=\sffamily\footnotesize] at (7.5,-1.5) {\qty{3}{cm}};
  \node[anchor=north, font=\footnotesize] at (7.5,-7) {\qty{9}{cm^2}};
  \bkiihcube{(21,0)}{3}{5}{bkTheme}
  \node[anchor=north, font=\sffamily\footnotesize] at (28.5,-1.5) {\qty{3}{cm}};
  \node[anchor=north, font=\footnotesize] at (31,-7) {\qty{27}{cm^3}};
\end{tikzpicture}
\end{minipage}

\unit{cm^2}（\R{平方}{へい|ほう}センチメートル）の2は、たて（cm）とよこ（cm）の2つの\R{長}{なが}さをかけたので、cmを2\R{回}{かい}かけている、という\R{意味}{い|み}です。\R{同}{おな}じように、\R{立方体}{りっ|ぽう|たい}の\R{大}{おお}きさは、たて・よこ・\R{高}{たか}さの3つの\R{長}{なが}さをかけるので、cmを3\R{回}{かい}かけた\unit{cm^3}で\R{表}{あらわ}します。1\R{辺}{ぺん}\qty{1}{cm}の\R{立方体}{りっ|ぽう|たい}の\R{大}{おお}きさを\qty{1}{cm^3}（1\R{立方}{りっ|ぽう}センチメートル）といい、\R{物}{もの}の\R{大}{おお}きさ（かさ）を\textbf{\R{体積}{たい|せき}}といいます。\bkidx{たいせき}{\R{体積}{たい|せき}}2\R{乗}{じょう}を\textbf{\R{平方}{へい|ほう}}、3\R{乗}{じょう}を\textbf{\R{立方}{りっ|ぽう}}ともいいます。

\noindent\begin{minipage}[c]{0.45\linewidth}
\hspace*{1\zw}\qty{1}{L}は、1\R{辺}{ぺん}\qty{10}{cm}の\R{立方体}{りっ|ぽう|たい}に\R{入}{はい}る\R{水}{みず}の\R{量}{りょう}と\R{決}{き}められています。その\R{体積}{たい|せき}は $10^3=10\times10\times10=1000$ で\qty{1000}{cm^3}。\qty{1}{L}は\qty{1000}{mL}でもあるので（1-9）、1000\R{等分}{とう|ぶん}した\qty{1}{mL}は、ちょうど\qty{1}{cm^3}です。
\end{minipage}\hfill
\begin{minipage}[c]{0.53\linewidth}\centering
\begin{tikzpicture}[x=1mm, y=1mm]
  \bkiihcube{(0,0)}{10}{1.7}{bkBlue}
  \node[anchor=north, font=\sffamily\footnotesize] at (8.5,-1.5) {\qty{10}{cm}};
  \node[anchor=north, font=\footnotesize, text=bkBlue!70!black] at (12,-7) {\qty{1000}{cm^3}＝\qty{1}{L}};
  \draw[bkGray!70, line width=0.6pt, dash pattern=on 2pt off 1.5pt] (28.5,-11) -- (28.5,28);
  \bkiihcube{(32,0)}{1}{19}{bkTheme}
  \bkiihcube{(34,2)}{1}{2.1}{bkBlue}
  \draw[bkInk!70, line width=0.5pt] (37.1,4.3) -- (40.5,8) node[anchor=south west, font=\sffamily\footnotesize, text=bkInk, inner sep=1pt] {\qty{1}{L}};
  \node[anchor=north, font=\sffamily\footnotesize] at (41.5,-1.5) {\qty{1}{m}};
  \node[anchor=north, font=\footnotesize, text=bkTheme!80!black] at (46,-7) {\qty{1}{m^3}＝\qty{1000}{L}};
  \bkkid{(67.2,0)}{28}{bkGreen}{bkHair}
  \node[anchor=north, font=\sffamily\footnotesize, text=bkGray] at (36,-13.5) {（\R{絵}{え}の\R{大}{おお}きさの\R{比}{ひ}はちがいます）};
\end{tikzpicture}
\end{minipage}

1\R{辺}{ぺん}\qty{1}{m}の\R{立方体}{りっ|ぽう|たい}の\R{体積}{たい|せき}は\qty{1}{m^3}（1\R{立方}{りっ|ぽう}メートル）です。これは\R{何}{なん}\unit{cm^3}でしょう。\qty{1}{m}は\qty{100}{cm}なので、1\R{辺}{ぺん}\qty{1}{cm}の\R{立方体}{りっ|ぽう|たい}が、1だんに $100\times100=10000$こならび、それが100だん\R{重}{かさ}なります。
\[ \qty{1}{m^3}=100\times100\times100\ \unit{cm^3}=\qty[group-digits=none]{1000000}{cm^3} \]
\qty{1000}{cm^3}が\qty{1}{L}なので、Lにするには1000でわります。
\[ 1000000\div1000=1000 \]
だから、\qty{1}{m^3}＝\qty{1000}{L}です。おふろの\R{水}{みず}（\R{約}{やく}\qty{200}{L}）の $1000\div200=5$ で、5はい\R{分}{ぶん}です。

\bksubsection{3}{\R{大}{おお}きな\R{数}{かず}は、10の\R{累乗}{るい|じょう}で}

\leavevmode\label{bk:2-7:juu}10を\R{何回}{なん|かい}もかけると、10、100、1000、……と、1のあとの0が1つずつふえます。10を4\R{回}{かい}かけた $10^4$ は10000で、1\R{万}{まん}です。

\noindent\begin{minipage}[c]{0.6\linewidth}
\begin{bktable}{colspec={Q[c,m,wd=12mm] Q[l,m,wd=36mm] Q[c,m,wd=14mm]}, rowsep=1.5pt, colsep=4pt}
 & \R{数字}{すう|じ} & \R{累乗}{るい|じょう} \\
1\R{万}{まん} & 1\,0000 & $10^4$ \\
1\R{億}{おく} & 1\,0000\,0000 & $10^8$ \\
1\R{兆}{ちょう} & 1\,0000\,0000\,0000 & $10^{12}$ \\
\end{bktable}
\end{minipage}\hfill
\begin{minipage}[c]{0.38\linewidth}
0が4つふえるごとに、\R{名前}{な|まえ}が\mbox{\R{万}{まん}→\R{億}{おく}→\R{兆}{ちょう}}と\R{変}{か}わります。\R{指数}{し|すう}は、1のあとにつく0の\R{数}{かず}です。
\end{minipage}

\R{光}{ひかり}は1\R{秒間}{びょう|かん}に、\R{約}{やく}300000000\,m\R{進}{すす}みます。0を\R{数}{かぞ}えるのはたいへんですね。300000000は、3のあとに0が8つ。100000000（1\R{億}{おく}）の3つ\R{分}{ぶん}なので、$300000000=3\times100000000=3\times10^8$ です。$3\times10^8$ と\R{書}{か}けば、ひと\R{目}{め}で「3\R{億}{おく}」とわかります。

\bkiihhead{3でない\R{数}{かず}のとき：\R{小数点}{しょう|すう|てん}を\R{動}{うご}かして\R{考}{かんが}える}
\R{地球}{ち|きゅう}から\R{月}{つき}までの\R{約}{やく}38\R{万}{まん}\,km（380000\,km）は、どう\R{書}{か}けばよいでしょう。10\R{倍}{ばい}すると\R{小数点}{しょう|すう|てん}が\R{右}{みぎ}へ1つ\R{動}{うご}きます（2-1）。3.8を10\R{倍}{ばい}、10\R{倍}{ばい}、……としていくと、

\begin{bkfigure}
\begin{tikzpicture}[x=1mm, y=1mm]
  \foreach \v [count=\i from 0] in {3.8, 38, 380, 3800, 38000, 380000} {
    \node[bk node, minimum width=16mm, minimum height=7mm, font=\small] (n\i) at (\i*23,0) {\v};}
  \foreach \i [evaluate=\i as \j using int(\i+1)] in {0,...,4} {
    \draw[bk link] (n\i) -- (n\j);
    \node[anchor=south, font=\sffamily\footnotesize, inner sep=0pt] at (\i*23+11.5,4.6) {×10};}
  \draw[bkTheme!80!black, line width=0.6pt, decorate, decoration={brace, amplitude=1.6mm, mirror}] (0,-5.5) -- (115,-5.5) node[midway, below=2.6mm, font=\sffamily\footnotesize, text=bkTheme!80!black] {10を5\R{回}{かい}かけた＝$\times10^5$};
\end{tikzpicture}
\bkfigcaption{3.8に10を5\R{回}{かい}かけると380000。だから $380000=3.8\times10^5$}
\end{bkfigure}

10を5\R{回}{かい}かけると380000になるので、$380000=3.8\times10^5$ です。\R{反対}{はん|たい}に、380000の\R{小数点}{しょう|すう|てん}を\R{左}{ひだり}へ5つ\R{動}{うご}かすと3.8になる、と\R{考}{かんが}えても\R{同}{おな}じです。\R{動}{うご}かした\R{数}{かず}が\R{指数}{し|すう}になります。

\noindent\textbf{なぜ×の\R{前}{まえ}を1\R{以上}{い|じょう}10より\R{小}{ちい}さくするの？}\quad $380000$ は $38\times10^4$ とも $0.38\times10^6$ とも\R{書}{か}けます。\R{書}{か}き\R{方}{かた}がいくつもあると、くらべにくいですね。そこで、×の\R{前}{まえ}を1\R{以上}{い|じょう}10より\R{小}{ちい}さい\R{数}{かず}にする\R{約束}{やく|そく}にします。すると、\R{指数}{し|すう}を\R{見}{み}るだけで\R{大}{おお}きさの\R{見当}{けん|とう}がつきます。$3.8\times10^5$（\R{月}{つき}まで）と $3\times10^8$（\R{光}{ひかり}が1\R{秒}{びょう}で\R{進}{すす}む\R{距離}{きょ|り}。ただしm）なら、\R{指数}{し|すう}の5と8がすぐ\R{目}{め}に\R{入}{はい}ります。

\R{理科}{り|か}では、\R{光}{ひかり}の\R{速}{はや}さを $3.0\times10^8$\,m/s と\R{書}{か}きます。ほんとうの\R{値}{あたい}は\R{約}{やく} $2.998\times10^8$\,m/s ですが、\R{有効数字}{ゆう|こう|すう|じ}2けた（2-6）に\R{四捨五入}{し|しゃ|ご|にゅう}して3.0です。3.0の0は、「\R{小数第}{しょう|すう|だい}1\R{位}{い}まで\R{考}{かんが}えて0」というしるしでしたね。

\bksubsection{4}{\R{指数}{し|すう}のきまり：かけ\R{算}{ざん}は\R{指数}{し|すう}をたす}

\leavevmode\label{bk:2-7:kimari}$10^2\times10^3$ はいくつでしょう。10を\R{全部}{ぜん|ぶ}で\R{何回}{なん|かい}かけているか、\R{数}{かぞ}えてみます。

\begin{bkfigure}
\begin{tikzpicture}[x=1mm, y=1mm, ten/.style={draw=#1!80!black, fill=#1!20, line width=0.7pt, rounded corners=1mm, minimum width=10mm, minimum height=7mm, inner sep=0pt, font=\sffamily\small}]
  \node[anchor=east, font=\small] at (-2,0) {$10^2\times10^3$};
  \node[font=\small] at (2.5,0) {＝};
  \foreach \i in {0,1} \node[ten=bkBlue] at (12+\i*12,0) {10};
  \node[font=\small] at (36,0) {×};
  \foreach \i in {0,1,2} \node[ten=bkOrange] at (48+\i*12,0) {10};
  \node[font=\small] at (80,0) {＝};
  \node[anchor=west, font=\small] at (84,0) {$10^5$};
  \draw[bkBlue!80!black, line width=0.6pt, decorate, decoration={brace, amplitude=1.6mm, mirror}] (7,-5) -- (29,-5) node[midway, below=2.4mm, font=\sffamily\footnotesize, text=bkBlue!80!black] {2\R{回}{かい}};
  \draw[bkOrange!80!black, line width=0.6pt, decorate, decoration={brace, amplitude=1.6mm, mirror}] (43,-5) -- (77,-5) node[midway, below=2.4mm, font=\sffamily\footnotesize, text=bkOrange!80!black] {3\R{回}{かい}};
  \node[anchor=west, font=\sffamily\footnotesize, text=bkInk] at (84,-8.6) {2＋3＝5\R{回}{かい}};
\end{tikzpicture}
\end{bkfigure}

10を2\R{回}{かい}と3\R{回}{かい}、あわせて5\R{回}{かい}かけています。\R{指数}{し|すう}は $2+3=5$ なので、\R{答}{こた}えは $10^5$ です。$100\times1000=100000$ で、たしかに0が5つです。2でも\R{同}{おな}じで、$2^3\times2^4=2^{3+4}=2^7=128$（$8\times16=128$）。\R{同}{おな}じ\R{数}{かず}の\R{累乗}{るい|じょう}どうしのかけ\R{算}{ざん}は、\textbf{\R{指数}{し|すう}をたせば}よいのです。

わり\R{算}{ざん}は、\R{分数}{ぶん|すう}の\R{形}{かたち}にして\R{約分}{やく|ぶん}します（2-4）。
\[ 10^5\div10^2=\frac{10\times10\times10\times\cancel{10}\times\cancel{10}}{\cancel{10}\times\cancel{10}}=10^3 \]
10が2つ\R{消}{き}えるので、\R{指数}{し|すう}は $5-2=3$。わり\R{算}{ざん}は、\textbf{\R{指数}{し|すう}をひけば}よいのです（$100000\div100=1000$）。

\par\addvspace{4pt}\noindent\bkhatenaref の\R{答}{こた}え\quad 42\R{回}{かい}\R{折}{お}った\R{紙}{かみ}は $2^{42}$\R{枚}{まい}。$42=10+10+10+10+2$ なので、\R{指数}{し|すう}のきまりから
\[ 2^{42}=2^{10}\times2^{10}\times2^{10}\times2^{10}\times2^2 \]
$2^{10}=1024$ は、ほぼ1000です。だから、10\R{回}{かい}\R{折}{お}るごとに、\R{厚}{あつ}さは\R{約}{やく}1000\R{倍}{ばい}になります。\R{電卓}{でん|たく}で、\qty{0.1}{mm}に1024をくり\R{返}{かえ}しかけていきましょう。\R{数}{かず}が\R{大}{おお}きくなったら、\R{単位}{たん|い}を\R{大}{おお}きくして\R{読}{よ}みやすくします。

\begin{bkcenter}
\begin{bktable}{colspec={Q[c,m,wd=11mm] Q[l,m,wd=42mm] Q[l,m,wd=60mm]}, rowsep=1.5pt, colsep=3pt, cells={font=\footnotesize}, row{1}={font=\sffamily\bfseries\footnotesize, bg=bkTheme!10}}
\R{折}{お}る & \R{計算}{けい|さん}（mm） & \R{単位}{たん|い}をなおす \\
10\R{回}{かい} & $0.1\times1024=102.4$ & \qty{1}{cm}＝\qty{10}{mm}なので÷10。\R{約}{やく}\qty{10}{cm} \\
20\R{回}{かい} & $102.4\times1024$ ＝\R{約}{やく}104858 & \qty{1}{m}＝\qty{1000}{mm}なので÷1000。\R{約}{やく}\qty{105}{m} \\
30\R{回}{かい} & \R{約}{やく}\qty{105}{m}の\R{約}{やく}1024\R{倍}{ばい} & \R{約}{やく}107374\,m。\qty{1}{km}＝\qty{1000}{m}なので÷1000。\R{約}{やく}\qty{107}{km} \\
40\R{回}{かい} & \R{約}{やく}\qty{107}{km}の\R{約}{やく}1024\R{倍}{ばい} & \R{約}{やく}109951\,km。\R{約}{やく}11\R{万}{まん}\,km \\
42\R{回}{かい} & 40\R{回}{かい}の $2^2=4$\R{倍}{ばい} & \R{約}{やく}439805\,km。\R{約}{やく}44\R{万}{まん}\,km \\
\end{bktable}
\bktabcaption{\R{単位}{たん|い}をなおすときは、1つ\R{上}{うえ}の\R{単位}{たん|い}の\R{中}{なか}に、\R{下}{した}の\R{単位}{たん|い}がいくつ\R{入}{はい}るかでわる。}
\end{bkcenter}

\R{数}{かず}の\R{大}{おお}きさを、\R{身近}{み|ぢか}なものにたとえてみると、こうなります。

\begin{bkcenter}
\begin{bktable}{colspec={Q[l,m,wd=15mm] Q[c,m,wd=23mm] Q[c,m,wd=19.5mm] Q[c,m,wd=18.5mm] Q[c,m,wd=23mm] Q[c,m,wd=20.5mm]}, rowsep=1.5pt, colsep=3pt, cells={font=\small}, row{1}={font=\sffamily\bfseries\small}, column{6}={bg=bkTheme!10}}
\R{折}{お}る\R{回数}{かい|すう} & 10\R{回}{かい} & 20\R{回}{かい} & 30\R{回}{かい} & 40\R{回}{かい} & 42\R{回}{かい} \\
\R{厚}{あつ}さ & \R{約}{やく}\qty{10}{cm} & \R{約}{やく}\qty{105}{m} & \R{約}{やく}\qty{107}{km} & \R{約}{やく}11\R{万}{まん}\,km & \R{約}{やく}44\R{万}{まん}\,km \\
たとえば & はがきのはば & ビルの30\R{階}{かい} & \R{宇宙}{う|ちゅう}に\R{入}{はい}る & \R{地球}{ち|きゅう}を\R{約}{やく}3\R{周}{しゅう} & \R{月}{つき}をこえる \\
\end{bktable}
\end{bkcenter}

（ビルは1\R{階}{かい}\R{約}{やく}\qty{3.5}{m}で、$105\div3.5=30$。\R{地球}{ち|きゅう}1\R{周}{しゅう}は\R{約}{やく}4\R{万}{まん}\,kmで、$11\div4=2.75$ で\R{約}{やく}3\R{周}{しゅう}。\R{宇宙}{う|ちゅう}は\R{地上}{ち|じょう}\R{約}{やく}\qty{100}{km}から。）

42\R{回}{かい}で\R{約}{やく}44\R{万}{まん}\,km。\R{月}{つき}までの\R{約}{やく}38\R{万}{まん}\,kmをこえるので、\R{答}{こた}えは「とどく」です（41\R{回}{かい}では\R{約}{やく}22\R{万}{まん}\,kmで、まだとどきません）。ただし、ほんとうの\R{紙}{かみ}は、7〜8\R{回}{かい}\R{折}{お}るのがやっとです。

\bksubsection{5}{0\R{乗}{じょう}とマイナスの\R{指数}{し|すう}}

\leavevmode\label{bk:2-7:zero}10の\R{累乗}{るい|じょう}を、\R{指数}{し|すう}を1つずつへらしてならべてみましょう。

\begin{bkfigure}
\begin{tikzpicture}[x=1mm, y=1mm]
  \foreach \e/\v/\c [count=\i from 0] in {3/1000/white, 2/100/white, 1/10/white, 0/？/bkTheme!14, -1/？/bkTheme!14, -2/？/bkTheme!14} {
    \node[bk node, minimum width=13mm, minimum height=9mm, inner xsep=1mm, fill=\c, font=\small] (b\i) at (\i*23.6,0) {$\v$};
    \node[anchor=south, font=\small, text=bkTheme!80!black] at (\i*23.6,6.5) {$10^{\e}$};}
  \foreach \i [evaluate=\i as \j using int(\i+1)] in {0,...,4} {
    \draw[bk link] (b\i) -- (b\j);
    \node[anchor=north, font=\sffamily\footnotesize, text=bkInk, inner sep=0pt] at (\i*23.6+11.8,-2.2) {÷10};}
\end{tikzpicture}
\end{bkfigure}

\R{指数}{し|すう}が1へると、\R{数}{かず}は10でわった\R{大}{おお}きさになっています。このきまりがつづくように、10でわって、その\R{先}{さき}を\R{決}{き}めます。
\[ 10^0=1,\qquad 10^{-1}=\frac{1}{10}=0.1,\qquad 10^{-2}=\frac{1}{10^2}=\frac{1}{100}=0.01 \]
2の\R{累乗}{るい|じょう}でも、$2^2=4$、$2^1=2$ から2でわっていくと、
\[ 2^0=1,\qquad 2^{-1}=\frac{1}{2},\qquad 2^{-2}=\frac{1}{2^2}=\frac{1}{4} \]
0\R{乗}{じょう}は、10でも2でも1。マイナスの\R{指数}{し|すう}は、「1をその\R{数}{かず}で\R{何回}{なん|かい}わるか」を\R{表}{あらわ}します。

\bkiihhead{わり\R{算}{ざん}のきまりからも、\R{同}{おな}じ\R{答}{こた}えになる}
「わり\R{算}{ざん}は\R{指数}{し|すう}をひく」のきまりで、$10^3\div10^3$ を\R{計算}{けい|さん}してみます。\R{同}{おな}じ\R{数}{かず}でわるので、\R{答}{こた}えは1です（$1000\div1000=1$）。\R{指数}{し|すう}をひくと $3-3=0$ で、$10^0$ です。2つの\R{答}{こた}えは\R{同}{おな}じはずなので、$10^0=1$ と\R{決}{き}めるのが\R{自然}{し|ぜん}です。

\R{同}{おな}じように、$10^2\div10^3$ は、$100\div1000=\dfrac{100}{1000}=\dfrac{1}{10}$。\R{指数}{し|すう}をひくと $2-3=-1$ で、$10^{-1}$ です。だから $10^{-1}=\dfrac{1}{10}$。\R{表}{ひょう}で\R{作}{つく}った\R{答}{こた}えと、ぴったり\R{合}{あ}いますね。

\R{紙}{かみ}の\R{枚数}{まい|すう}でも\R{確}{たし}かめられます。0\R{回}{かい}\R{折}{お}った（まだ\R{折}{お}っていない）\R{紙}{かみ}は1\R{枚}{まい}で、$2^0=1$ です（\bkyattemiyouref の\R{表}{ひょう}の0\R{回}{かい}）。

こう\R{決}{き}めると、\R{指数}{し|すう}のきまりがそのまま\R{使}{つか}えます。$10^3\times10^{-1}=10^{3+(-1)}=10^2$（$3+(-1)=2$ は2-2のたし\R{算}{ざん}）は、$1000\times\dfrac{1}{10}=100$ とぴったり\R{合}{あ}います。

\noindent\textbf{$10^{-1}$ はマイナスの\R{数}{かず}ではない}\quad $10^{-1}=0.1$ は、0より\R{大}{おお}きい\R{数}{かず}です。\R{指数}{し|すう}のマイナスは「わる」ことを\R{表}{あらわ}していて、\R{数}{かず}がマイナスになるわけではありません。$10^{-3}=\dfrac{1}{1000}=0.001$ は、\qty{1}{m}の1000\R{分}{ぶん}の1、つまり\qty{1}{mm}の\R{大}{おお}きさです（2-10でまた\R{使}{つか}います）。


\begin{motto}[2]
\par\addvspace{3pt}\textbf{\R{累乗}{るい|じょう}の\R{累乗}{るい|じょう}}\par\noindent $10^2$ を3\R{回}{かい}かけると、\R{指数}{し|すう}は $2+2+2=6$ なので、$(10^2)^3=10^6$ です。$\qty{1}{m}=10^2$\,cm なので、\qty{1}{m^3}は $(10^2)^3=10^6$\,\unit{cm^3}。\R{数学}{すう|がく}の\R{道具}{どう|ぐ}2で\R{求}{もと}めた1000000\,\unit{cm^3}（0が6つ）と\R{合}{あ}います。
\end{motto}

\begin{taisetsu}
\R{同}{おな}じ\R{数}{かず}を3\R{回}{かい}かけた\R{数}{かず}は、$2\times2\times2=2^3$ のように\R{書}{か}く。\\
$10^2\times10^3=10^{2+3}=10^5$\quad（かけ\R{算}{ざん}は\R{指数}{し|すう}をたす）\\
$10^5\div10^2=10^{5-2}=10^3$\quad（わり\R{算}{ざん}は\R{指数}{し|すう}をひく）\\
$10^0=1$、\ $10^{-1}=\dfrac{1}{10}$、\ $10^{-2}=\dfrac{1}{10^2}=\dfrac{1}{100}$
\end{taisetsu}

\begin{kotoba}
\small
\bkword{\R{累乗}{るい|じょう}}{\R{同}{おな}じ\R{数}{かず}をいくつかかけたもの。$2^3$ は「2の3\R{乗}{じょう}」。}
\bkword{\R{指数}{し|すう}}{\R{右}{みぎ}かたの\R{小}{ちい}さな\R{数}{かず}。\R{何回}{なん|かい}かけるかを\R{表}{あらわ}す。}
\bkword{\R{平方}{へい|ほう}・\R{立方}{りっ|ぽう}}{2\R{乗}{じょう}と3\R{乗}{じょう}のこと。\unit{cm^2}の2、\unit{cm^3}の3。}
\bkword{\R{体積}{たい|せき}}{\R{物}{もの}の\R{大}{おお}きさ（かさ）。$\qty{1}{mL}=\qty{1}{cm^3}$、$\qty{1}{L}=\qty{1000}{cm^3}$}
\end{kotoba}

\bkkigocards{\bkkigocard{$2^3$}{2の3\R{乗}{じょう}}{2を3\R{回}{かい}かけた\R{数}{かず}。$2\times2\times2=8$}{2 to the 3rd power}{2-7}}
{\bkkigocard{\sffamily\unit{cm^3}}{\R{立方}{りっ|ぽう}センチメートル}{\R{体積}{たい|せき}の\R{単位}{たん|い}。$\qty{1}{L}=\qty{1000}{cm^3}$}{cubic centimeter}{2-7}}

%% ---------------- 例題 ----------------
\bksection{reidai}{\R{例題}{れい|だい}でたしかめよう}

\begin{reidai}[\R{累乗}{るい|じょう}を\R{計算}{けい|さん}する]
（1）$3^4$\quad（2）$1^5$\quad（3）$10^5$ を\R{計算}{けい|さん}しましょう。\par
\kangaekata \R{指数}{し|すう}の\R{数}{かず}だけ、\R{同}{おな}じ\R{数}{かず}をかけます。\R{前}{まえ}から\R{順}{じゅん}に1つずつかけます。\par
\shiki （1）$3^4=3\times3\times3\times3=9\times3\times3=27\times3=81$\par
\hspace*{2.6\zw}（2）$1^5=1\times1\times1\times1\times1=1$\quad（1は\R{何回}{なん|かい}かけても1）\par
\hspace*{2.6\zw}（3）$10^5=10\times10\times10\times10\times10=100000$\par
\kotae （1）81\quad（2）1\quad（3）100000
\end{reidai}

\begin{jibun}
（1）$4^3$\quad（2）$1^8$\quad（3）$10^6$ を\R{計算}{けい|さん}しましょう。
\bkjans{（1）\nobreak 64\quad（2）\nobreak 1\quad（3）\nobreak 1000000}
\end{jibun}

\begin{reidai}[\R{時間}{じ|かん}を\R{分}{ふん}になおして、\R{倍}{ばい}になる\R{回数}{かい|すう}を\R{求}{もと}める]
20\R{分}{ぷん}ごとに2つに\R{分}{わ}かれてふえる\R{細菌}{さい|きん}があります。1こが2\R{時間後}{じ|かん|ご}には\R{何}{なん}こになりますか。\par
\kangaekata 「20\R{分}{ぷん}ごと」なので、\R{時間}{じ|かん}を\R{分}{ふん}にそろえます。1\R{時間}{じ|かん}は60\R{分}{ぷん}です。2\R{時間}{じ|かん}の\R{中}{なか}に20\R{分}{ぷん}がいくつあるかが、2\R{倍}{ばい}になる\R{回数}{かい|すう}です。\par
\shiki ① 2\R{時間}{じ|かん}は、60\R{分}{ぷん}の2つ\R{分}{ぶん}。$60\times2=120$ で120\R{分}{ぷん}\par
\hspace*{2.6\zw}② 2\R{倍}{ばい}になる\R{回数}{かい|すう}は $120\div20=6$（\R{回}{かい}）\par
\hspace*{2.6\zw}③ 1こを6\R{回}{かい}2\R{倍}{ばい}にするので、$1\times2^6=2\times2\times2\times2\times2\times2=64$\par
\kotae 64こ
\end{reidai}

\begin{jibun}
\bkreidairef{2}の\R{細菌}{さい|きん}は、1こが3\R{時間後}{じ|かん|ご}には\R{何}{なん}こになりますか。
\bkjans{3\R{時間}{じ|かん}は $60\times3=180$ で180\R{分}{ぷん}。$180\div20=9$\R{回}{かい}なので $2^9=512$ で512こ}
\end{jibun}

\begin{reidai}[\R{立方体}{りっ|ぽう|たい}の\R{体積}{たい|せき}とL]
1\R{辺}{ぺん}\qty{20}{cm}の\R{立方体}{りっ|ぽう|たい}の\R{箱}{はこ}に\R{入}{はい}る\R{水}{みず}は、\R{何}{なん}\unit{cm^3}、\R{何}{なん}Lですか。\par
\kangaekata 1\R{辺}{ぺん}\qty{1}{cm}の\R{立方体}{りっ|ぽう|たい}が、1だんに $20\times20$こ、それが20だん。\qty{1}{L}は\qty{1000}{cm^3}なので、Lにするには1000でわります。\par
\shiki ① $20^3=20\times20\times20=400\times20=8000$\par
\hspace*{2.6\zw}② $8000\div1000=8$\par
\kotae \qty{8000}{cm^3}、\qty{8}{L}
\end{reidai}

\begin{jibun}
1\R{辺}{ぺん}\qty{30}{cm}の\R{立方体}{りっ|ぽう|たい}の\R{水}{みず}そうに\R{入}{はい}る\R{水}{みず}は、\R{何}{なん}\unit{cm^3}、\R{何}{なん}Lですか。
\bkjans{$30^3=27000$ で\qty[group-digits=none]{27000}{cm^3}、\qty{27}{L}}
\end{jibun}

\begin{reidai}[\R{大}{おお}きな\R{数}{かず}を10の\R{累乗}{るい|じょう}で]
（1）5\R{億}{おく}、（2）1\R{億}{おく}5000\R{万}{まん}を、□$\times10^{\text{△}}$ の\R{形}{かたち}で\R{書}{か}きましょう。（3）$4\times10^6$ を\R{数字}{すう|じ}で\R{書}{か}きましょう。\par
\kangaekata 1\R{億}{おく}は100000000（1のあとに0が8つ）なので $10^8$。□は1\R{以上}{い|じょう}10より\R{小}{ちい}さい\R{数}{かず}にします。\par
\shiki （1）5\R{億}{おく}は1\R{億}{おく}の5つ\R{分}{ぶん}。$500000000=5\times100000000=5\times10^8$\par
\hspace*{2.6\zw}（2）150000000の\R{小数点}{しょう|すう|てん}を\R{左}{ひだり}へ8つ\R{動}{うご}かすと1.5\par
\hspace*{2.6\zw}\phantom{（2）}だから $150000000=1.5\times100000000=1.5\times10^8$\par
\hspace*{2.6\zw}（3）$10^6=1000000$ なので、$4\times10^6=4\times1000000=4000000$\par
\kotae （1）$5\times10^8$\quad（2）$1.5\times10^8$\quad（3）4000000（400\R{万}{まん}）
\end{reidai}

\begin{jibun}
（1）7\R{兆}{ちょう}、（2）1\R{億}{おく}2000\R{万}{まん}を、□$\times10^{\text{△}}$ の\R{形}{かたち}で\R{書}{か}きましょう。（3）$2.5\times10^5$ を\R{数字}{すう|じ}で\R{書}{か}きましょう。
\bkjans{（1）\nobreak $7\times10^{12}$\quad（2）\nobreak $1.2\times10^8$\quad（3）\nobreak 250000（25\R{万}{まん}）}
\end{jibun}

\begin{reidai}[\R{指数}{し|すう}のきまりを\R{使}{つか}う]
（1）$10^4\times10^3$\quad（2）$2^7\div2^3$\quad（3）$10^5\times10^{-2}$ を\R{累乗}{るい|じょう}で、（4）$10^{-3}$ を\R{小数}{しょう|すう}で\R{表}{あらわ}しましょう。\par
\kangaekata かけ\R{算}{ざん}は\R{指数}{し|すう}をたす。わり\R{算}{ざん}は\R{指数}{し|すう}をひく。マイナスの\R{指数}{し|すう}は、1をその\R{数}{かず}で\R{何回}{なん|かい}わるか。\par
\shiki （1）$10^{4+3}=10^7$\qquad（2）$2^{7-3}=2^4$\quad（$128\div8=16$）\par
\hspace*{2.6\zw}（3）$10^{5+(-2)}=10^3$\qquad（4）$10^{-3}=\dfrac{1}{10^3}=\dfrac{1}{1000}=0.001$\par
\kotae （1）$10^7$\quad（2）$2^4$（16）\quad（3）$10^3$（1000）\quad（4）0.001
\end{reidai}

\begin{jibun}
（1）$10^2\times10^6$\quad（2）$10^8\div10^5$\quad（3）$10^{-1}\times10^{-2}$ を\R{累乗}{るい|じょう}で、（4）$2^{-3}$ を\R{分数}{ぶん|すう}で\R{表}{あらわ}しましょう。
\bkjans{（1）\nobreak $10^8$\quad（2）\nobreak $10^3$\quad（3）\nobreak $10^{-3}$\quad（4）\nobreak $\dfrac{1}{8}$}
\end{jibun}

\begin{column}{\R{紙}{かみ}は、ほんとうは\R{何回}{なん|かい}\R{折}{お}れる？}
コピー\R{用紙}{よう|し}は、7〜8\R{回}{かい}\R{折}{お}るのがやっとです。\R{折}{お}るたびに\R{厚}{あつ}さは2\R{倍}{ばい}、\R{広}{ひろ}さは\R{半分}{はん|ぶん}になるからです。2002\R{年}{ねん}、アメリカの\R{高校生}{こう|こう|せい}が、\R{長}{なが}さ\R{約}{やく}\qty{1.2}{km}のトイレットペーパーを12\R{回}{かい}\R{折}{お}ることに\R{成功}{せい|こう}しました。$2^{12}=4096$\R{枚}{まい}の\R{重}{かさ}なりです。
\end{column}

\begin{reidai}[はじめて□をこえるのは\R{何回}{なん|かい}め？]
\R{厚}{あつ}さ\qty{0.1}{mm}の\R{紙}{かみ}を、\R{何回}{なん|かい}でも\R{折}{お}れるとします。\R{厚}{あつ}さがはじめて\qty{10}{cm}をこえるのは、\R{何回}{なん|かい}\R{折}{お}ったときですか。\par
\kangaekata \R{単位}{たん|い}をmmにそろえます。\qty{1}{cm}＝\qty{10}{mm}なので、\qty{10}{cm}＝\qty{100}{mm}。\R{厚}{あつ}さは「$0.1\times$（2の\R{累乗}{るい|じょう}）」なので、2の\R{累乗}{るい|じょう}の\R{表}{ひょう}を\R{順}{じゅん}に\R{見}{み}て、100をこえる\R{所}{ところ}をさがします。\par
\begin{bkcenter}
\begin{bktable}{colspec={Q[l,m,wd=26mm] *{4}{Q[c,m,wd=15mm]}}, rowsep=1pt, colsep=3pt, cells={font=\small}, column{1}={font=\sffamily\small}, cell{2}{5}={bg=bkOrange!22}}
\R{折}{お}った\R{回数}{かい|すう} & 7 & 8 & 9 & 10 \\
\R{厚}{あつ}さ（mm） & 12.8 & 25.6 & 51.2 & 102.4 \\
\end{bktable}
\end{bkcenter}
\shiki ① 9\R{回}{かい}：$0.1\times2^9=0.1\times512=51.2$（mm）\par
\hspace*{2.6\zw}\phantom{①}\ まだ100にとどかない\par
\hspace*{2.6\zw}② 10\R{回}{かい}：$0.1\times2^{10}=0.1\times1024=102.4$（mm）\par
\hspace*{2.6\zw}\phantom{②}\ 100をこえた\par
\kotae 10\R{回}{かい}（\R{約}{やく}\qty{10.2}{cm}）
\end{reidai}

\begin{jibun}
1この\R{細菌}{さい|きん}が20\R{分}{ぷん}ごとに2\R{倍}{ばい}にふえます。はじめて100こをこえるのは\R{何}{なん}\R{分後}{ぷん|ご}ですか。また、それは\R{何時間何分後}{なん|じ|かん|なん|ぷん|ご}ですか。
\bkjans{$2^6=64$、$2^7=128$ なので7\R{回}{かい}。$20\times7=140$ で140\R{分後}{ぷん|ご}。$140=60\times2+20$ なので2\R{時間}{じ|かん}20\R{分後}{ぷん|ご}}
\end{jibun}



%% ---------------- まちがえやすいところ・物理工学 ----------------
\bksection{tsumazuki}{まちがえやすいところをチェック}

\begin{tsumazuki}{\R{指数}{し|すう}を、かける\R{数}{かず}にしてしまう}
\bkng{$2^3=2\times3=6$}
\bkok{$2^3=2\times2\times2=8$}
\R{指数}{し|すう}の3は、2を「3\R{回}{かい}かける」という\R{意味}{い|み}です。
\end{tsumazuki}

\begin{tsumazuki}{かけ\R{算}{ざん}で、\R{指数}{し|すう}どうしをかけてしまう}
\bkng{$10^2\times10^3=10^{2\times3}=10^6$}
\bkok{$10^2\times10^3=10^{2+3}=10^5$}
10を2\R{回}{かい}と3\R{回}{かい}、あわせて5\R{回}{かい}かけています（0は5つ）。
\end{tsumazuki}

\begin{tsumazuki}{マイナスの\R{指数}{し|すう}を、マイナスの\R{数}{かず}にしてしまう}
\bkng{$10^{-1}=-10$}
\bkok{$10^{-1}=\dfrac{1}{10}=0.1$}
\R{指数}{し|すう}のマイナスは「1を10で\R{何回}{なん|かい}わるか」です。\R{答}{こた}えは0より\R{大}{おお}きい\R{数}{かず}です。
\end{tsumazuki}

\begin{tsumazuki}{$2^3$ と $3^2$ を\R{同}{おな}じだと\R{思}{おも}う}
\bkng{$2^3=3^2$}
\bkok{$2^3=2\times2\times2=8$、\ $3^2=3\times3=9$}
\R{下}{した}の\R{大}{おお}きい\R{数}{かず}が「かける\R{数}{かず}」、\R{右}{みぎ}かたの\R{数}{かず}が「かける\R{回数}{かい|すう}」です。\R{入}{い}れかえると\R{別}{べつ}の\R{数}{かず}になります。
\end{tsumazuki}



\begin{butsuri}
\noindent\begin{minipage}[c]{0.73\linewidth}
\par\addvspace{3pt}\textbf{スイッチとメモリー}\par\noindent スイッチ1つは、\R{入}{いり}（1）か\R{切}{きり}（0）の2\R{通}{とお}り。3つなら $2^3=8$\R{通}{とお}り、10こなら $2^{10}=1024$\R{通}{とお}り。コンピュータのメモリーは、0か1をおぼえる\R{小}{ちい}さなスイッチの\R{集}{あつ}まりです（2-17）。メモリーの\R{大}{おお}きさでは、$2^{10}=1024$ を「1K（キロ）」とよぶことがあります（ふつうの\qty{1}{km}のkは1000）。
\end{minipage}\hfill
\begin{minipage}[c]{0.25\linewidth}\centering
\begin{tikzpicture}[x=1mm, y=1mm]
  \foreach \b [count=\i from 0] in {1,0,1} \bkiihswitch{(\i*10,0)}{\b}
  \node[anchor=north, font=\sffamily\footnotesize, text=bkInk] at (13.5,-5) {3つで $2^3=8$\R{通}{とお}り};
\end{tikzpicture}
\end{minipage}

\vspace{8pt}
\noindent\begin{minipage}[c]{0.73\linewidth}
\par\addvspace{3pt}\textbf{\R{細菌}{さい|きん}と\R{冷蔵庫}{れい|ぞう|こ}}\par\noindent \R{食中毒}{しょく|ちゅう|どく}の\R{原因}{げん|いん}になる\R{細菌}{さい|きん}には、あたたかい\R{所}{ところ}で20\R{分}{ぷん}ほどで2\R{倍}{ばい}にふえるものがあります。10\R{時間}{じ|かん}は $60\times10=600$ で600\R{分}{ぷん}、$600\div20=30$ で30\R{回}{かい}2\R{倍}{ばい}になるので、$2^{30}=1073741824$、\R{約}{やく}10\R{億}{おく}\R{倍}{ばい}です。\R{冷蔵庫}{れい|ぞう|こ}で\R{冷}{ひ}やすと、2\R{倍}{ばい}になるまでの\R{時間}{じ|かん}がとても\R{長}{なが}くなります。
\end{minipage}\hfill
\begin{minipage}[c]{0.25\linewidth}\centering
\begin{tikzpicture}[x=0.85mm, y=1mm]
  \bkiihcell{(0,0)}{3.2}
  \draw[bk arrow, bkGray, line width=0.9pt] (4.5,0) -- (9.5,0);
  \bkiihcell{(14,3.2)}{3.2}\bkiihcell{(14,-3.2)}{3.2}
  \draw[bk arrow, bkGray, line width=0.9pt] (18.5,0) -- (23.5,0);
  \foreach \yy in {-7.5,-2.5,2.5,7.5} \bkiihcell{(28,\yy)}{2.6}
  \node[anchor=north, font=\sffamily\footnotesize] at (14,-10.5) {20\R{分}{ぷん}ごとに2\R{倍}{ばい}};
\end{tikzpicture}
\end{minipage}
\end{butsuri}


\clearpage
%% ---------------- 章末問題 ----------------
\bkendproblems
\bkproblemgroup{\R{基本}{き|ほん}}{1}

\bkproblem \R{次}{つぎ}の\R{計算}{けい|さん}をしましょう。\quad（1）$2^4$\quad（2）$5^3$\quad（3）$10^5$\quad（4）$1^{10}$
\bkpans{（1）\nobreak 16\quad（2）\nobreak 125\quad（3）\nobreak 100000\quad（4）\nobreak 1}

\bkproblem 1\R{辺}{ぺん}\qty{10}{cm}の\R{立方体}{りっ|ぽう|たい}の\R{体積}{たい|せき}は\R{何}{なん}\unit{cm^3}ですか。それは\R{何}{なん}Lですか。また、\qty{1}{m^3}は\R{何}{なん}Lですか。
\bkpans{$10^3=1000$ で\qty{1000}{cm^3}、\qty{1}{L}。$\qty{1}{m^3}=\qty[group-digits=none]{1000000}{cm^3}=\qty{1000}{L}$}

\bkproblem \R{次}{つぎ}の\R{数}{かず}を、□$\times10^{\text{△}}$ の\R{形}{かたち}で\R{書}{か}きましょう（□は1\R{以上}{い|じょう}10より\R{小}{ちい}さい\R{数}{かず}）。\par
（1）3\R{億}{おく}\quad（2）6\R{兆}{ちょう}\quad（3）25\R{万}{まん}\quad（4）\R{地球}{ち|きゅう}から\R{太陽}{たい|よう}までの\R{約}{やく}150000000\,km
\bkpans{（1）\nobreak $3\times10^8$\quad（2）\nobreak $6\times10^{12}$\quad（3）\nobreak $2.5\times10^5$\quad（4）\nobreak $1.5\times10^8$\,km}

\bkproblem \R{次}{つぎ}の\R{計算}{けい|さん}をして、\R{答}{こた}えを\R{累乗}{るい|じょう}の\R{形}{かたち}で\R{書}{か}きましょう。\par
（1）$10^3\times10^5$\quad（2）$2^4\times2^3$\quad（3）$10^9\div10^4$\quad（4）$2^{10}\div2^{2}$
\bkpans{（1）\nobreak $10^8$\quad（2）\nobreak $2^7$\quad（3）\nobreak $10^5$\quad（4）\nobreak $2^8$}

\bkproblem \R{次}{つぎ}の\R{数}{かず}を、\R{整数}{せい|すう}か\R{分数}{ぶん|すう}か\R{小数}{しょう|すう}で\R{表}{あらわ}しましょう。\par
（1）$10^0$\quad（2）$10^{-2}$（\R{小数}{しょう|すう}で）\quad（3）$2^{-1}$（\R{分数}{ぶん|すう}で）\quad（4）$10^4\times10^{-4}$
\bkpans{（1）\nobreak 1\quad（2）\nobreak 0.01\quad（3）\nobreak $\dfrac{1}{2}$\quad（4）\nobreak $10^0=1$}

\bkproblem \R{次}{つぎ}の\R{数}{かず}を、ふつうの\R{数字}{すう|じ}で\R{書}{か}きましょう。\par
（1）\nobreak $10^6$\quad（2）\nobreak $2.5\times10^4$\quad（3）\nobreak $7\times10^8$
\bkpans{（1）\nobreak 1000000\quad（2）\nobreak 25000\quad（3）\nobreak 700000000（7\R{億}{おく}）}

\bkproblemgroup{チャレンジ}{2}

\bkproblem \R{厚}{あつ}さ\qty{0.1}{mm}の\R{紙}{かみ}を、\R{何回}{なん|かい}でも\R{折}{お}れるとします。\R{厚}{あつ}さが\R{初}{はじ}めて\qty{1}{m}（\qty{1000}{mm}）をこえるのは、\R{何回}{なん|かい}\R{折}{お}ったときですか。
\bkpans{\qty{1}{m}＝\qty{1000}{mm}。$0.1\times2^{13}=819.2$\,mm、$0.1\times2^{14}=1638.4$\,mm なので、14\R{回}{かい}}

\bkproblem \R{細菌}{さい|きん}が20\R{分}{ぷん}ごとに2\R{倍}{ばい}にふえます。1こが1000こをこえるのは、\R{何分後}{なん|ぷん|ご}ですか。\R{表}{ひょう}のつづきをノートに\R{書}{か}いて\R{考}{かんが}えましょう。
\begin{bkcenter}
\begin{bktable}{colspec={Q[l,m,wd=24mm] *{6}{Q[c,m,wd=12mm]}}, rowsep=1.5pt, colsep=3pt, row{1}={bg=white, font=\small}, column{1}={font=\sffamily\small}}
\R{時間}{じ|かん}（\R{分}{ふん}） & 0 & 20 & 40 & 60 & 80 & …… \\
\R{細菌}{さい|きん}の\R{数}{かず}（こ） & 1 & 2 & 4 & 8 & ？ & …… \\
\end{bktable}
\end{bkcenter}
\bkpans{$2^9=512$、$2^{10}=1024$ なので、10\R{回}{かい}ふえたとき。$20\times10=200$ で200\R{分後}{ぷん|ご}。$200=60\times3+20$ なので3\R{時間}{じ|かん}20\R{分後}{ぷん|ご}}

\bkproblem $10^3\times10^{-3}$ を\R{指数}{し|すう}のきまりで\R{計算}{けい|さん}すると、$10^{\text{□}}$ になります。□にあてはまる\R{数}{かず}と、その\R{答}{こた}えが1になることを、\R{数}{かず}で\R{確}{たし}かめましょう。
\bkpans{$10^{3+(-3)}=10^0$ で、□は0。$1000\times\dfrac{1}{1000}=1$ で、$10^0=1$ と\R{合}{あ}う}

\begin{column}{\R{兆}{ちょう}の\R{次}{つぎ}は？}
\R{兆}{ちょう}の1\R{万}{まん}\R{倍}{ばい}が1\R{京}{けい}（$10^{16}$）、その1\R{万}{まん}\R{倍}{ばい}が1\R{垓}{がい}（$10^{20}$）。\R{名前}{な|まえ}は0が4つふえるごとにつづき、いちばん\R{大}{おお}きな「\R{無量大数}{む|りょう|たい|すう}」は $10^{68}$ です。2012\R{年}{ねん}に\R{完成}{かん|せい}した\R{日本}{に|ほん}のスーパーコンピュータ「\R{京}{けい}」は、1\R{秒間}{びょう|かん}に1\R{京}{けい}\R{回}{かい}（$10^{16}$\R{回}{かい}）の\R{計算}{けい|さん}ができることから\R{名}{な}づけられました。
\end{column}

\begin{furikaeri}
\begin{checklist}
\item \R{同}{おな}じ\R{数}{かず}を\R{何回}{なん|かい}もかけた\R{数}{かず}を、\R{累乗}{るい|じょう}で\R{書}{か}いて\R{計算}{けい|さん}できる。\hfill{\footnotesize\sffamily\color{bkTheme}→\pageref{bk:2-7:ruijo}ページ}
\item \R{体積}{たい|せき}の\R{単位}{たん|い}や\R{大}{おお}きな\R{数}{かず}を、10の\R{累乗}{るい|じょう}で\R{表}{あらわ}せる。\hfill{\footnotesize\sffamily\color{bkTheme}→\pageref{bk:2-7:juu}ページ}
\item \R{指数}{し|すう}のきまり、0\R{乗}{じょう}、マイナスの\R{指数}{し|すう}がわかる。\hfill{\footnotesize\sffamily\color{bkTheme}→\pageref{bk:2-7:kimari}ページ}
\end{checklist}
\end{furikaeri}

\bkfukidashi{mina}{\R{指数}{し|すう}を\R{使}{つか}うと、\R{月}{つき}までの\R{距離}{きょ|り}も、とても\R{小}{ちい}さな\R{数}{かず}も、\R{短}{みじか}く\R{書}{か}けるね。\R{次}{つぎ}の2-8では、2\R{乗}{じょう}を「\R{逆}{ぎゃく}から」\R{考}{かんが}えるよ。}
